\documentclass[11pt]{article}

% arXiv cs.MS (cross-list cs.LO) submission.
\usepackage[margin=1.1in]{geometry}
\usepackage[T1]{fontenc}
\usepackage[utf8]{inputenc}
\usepackage{amsmath,amssymb,amsthm}
\usepackage{booktabs}
\usepackage{microtype}
\usepackage[hidelinks]{hyperref}
\usepackage{url}
\usepackage{listings}
\usepackage{xcolor}

\lstdefinelanguage{lean}{
  keywords={structure, inductive, theorem, def, lemma, instance, namespace,
            end, import, where, by, fun, match, with, deriving},
  sensitive=true,
  comment=[l]{--},
  morestring=[b]",
}
\lstset{
  language=lean,
  basicstyle=\small\ttfamily,
  keywordstyle=\bfseries,
  commentstyle=\itshape\color{gray},
  columns=fullflexible,
  breaklines=true,
  frame=single,
  framerule=0.3pt,
  xleftmargin=1em,
}

\newtheorem{theorem}{Theorem}
\newtheorem{definition}{Definition}

% Release-mode guard: release builds hard-fail on missing generated
% fragments; draft builds show placeholders. Volatile standing claims live
% only in generated/ (see generated/FRAGMENT_CONTRACT.md).
\newif\ifreleasebuild
\releasebuildtrue
% GENERATED FILE — DO NOT EDIT BY HAND
% Source: release-manifest.json + standing.env via the quadrature graph
% Generator: ggen (quadrature-pack) — ggen renders; Lean admits; mfact certifies.
% Standing: GENERATED_FROM_CERTIFIED_MANIFEST
\newcommand{\ReleaseTag}{v26.7.6-procint-certified}
\newcommand{\ReleaseHash}{ff62f3e569af99977c2e7f5eb0b96efd835d00d530a39b442551654ae385605e}
\newcommand{\ReleaseRun}{e1ddbc6}
\newcommand{\SorryCount}{0}
\newcommand{\AdmitCount}{0}
\newcommand{\KernelTheorems}{145}
\newcommand{\TotalDecls}{309}
\newcommand{\StatedCount}{2}
\newcommand{\AxiomAuditStatus}{PASS}
\newcommand{\FixtureStatus}{PASS}
\newcommand{\CertifiedStatus}{PASS}
\newcommand{\QuadratureStatus}{PASS}
\newcommand{\CrownJewelStatus}{STATED}


\title{Receipted Mathematical Manufacturing:\\A Lean 4 Pipeline for Process Intelligence\\
\large mfact manufactures procint}
\author{Sean Chatman}
\date{}

\begin{document}
\maketitle

\begin{abstract}
We present \textsf{mfact}, a Lean~4 framework for \emph{mathematical
manufacturing}: the production of formal mathematical artifacts whose
standing is determined not by authorship, generation, or plausibility, but
by typed admission, refusal, audit, and receipt. The governing law is
\(R_B \vdash A = \mu(O^*_B)\): inside a declared boundary \(B\), raw
observations and generated candidates become artifacts with standing only
after they are admitted as \(O^*_B\), transformed by a lawful manufacturing
function \(\mu\), and accompanied by replayable receipts \(R_B\). This
separates contact from consequence. Humans, templates, and LLM agents may
produce candidates; only the Lean kernel, Lake build, axiom audit, sorry
census, negative fixtures, and release manifest can confer standing.

As the first manufactured artifact we introduce \textsf{procint}, a Lean~4
library for process intelligence: Petri nets with finitely supported
markings, workflow nets with soundness split into independent clauses,
event logs and alignment-based conformance, and object-centric event logs
following OCEL~2.0. The library is manufactured from a declared obligation
graph: an RDF/Turtle catalog of structures, theorem surfaces, module
dependencies, fixtures, citations, and registry surfaces is projected
through SPARQL-in-Tera templates by \textsf{ggen}, then admitted or
refused by Lean/Lake. Every failed admission is a typed refusal rather
than a silent downgrade; every admitted declaration is recorded with its
hash, axiom footprint, fixture status, and proven/stated standing. The
contribution is not merely that \textsf{procint} exists; it is that
\textsf{procint} is the first product of \textsf{mfact}.

The result is not a claim that arbitrary generated mathematics is
automatically correct. Rice-style undecidability is respected rather than
ignored: \textsf{mfact} replaces the impossible question ``is this
arbitrary generated artifact semantically correct?'' with the bounded,
replayable question ``did this candidate pass the declared admission gate
inside boundary \(B\)?'' The evaluation section of this paper is generated
from the release manifest rather than written by hand, and a
kernel-checked \emph{Standing Quadrature} witness closes the cross-product
between the declaration catalog, the admitted corpus, the release
manifest, the process evidence, and this paper's claims. We report exactly
what was proven, what remains stated, what was refused, and what the
release boundary excludes.
\end{abstract}

\section{Introduction}

Formalization today often resembles hand-fitted manufacturing: every
theorem, lemma, import, and repair is shaped locally by expert craft. The
result can be excellent, but it is not yet industrial. Mathematical
manufacturing asks a different question: what must be standardized so that
formal domains can be produced as interchangeable parts? In
\textsf{mfact}, the interchangeable part is not merely a Lean declaration.
It is a declaration plus provenance, admission status, axiom footprint,
refusal mode, hash, manifest position, and replay path. A theorem without
this surrounding receipt is hand-fitted metal; an admitted declaration
with its release evidence is an interchangeable mathematical part. The
manufactured object is not syntax --- it is \emph{standing}.

Recent work uses LLMs to draft formal content, which raises a question
the kernel alone does not settle: what standing does generated material
have? Kernel acceptance establishes that a term inhabits its stated type;
it does not establish that the stated type is the intended one, that the
statement is not vacuously weakened, or that the artifact as a whole is a
faithful rendering of a target theory. Meek et
al.~\cite{meek2026formalizing} document these failure modes in detail:
compilation masks incomplete multi-part statements, added weakening
hypotheses, and parameter restrictions. Our position is compressed by one
rule: \emph{contact is not consequence}. LLM contact with a theorem is
not consequence. A model may propose a proof term, a definition, or a
module, but the proposal has no standing until admitted. \textsf{mfact}
formalizes this separation: contact produces a candidate; admission
produces consequence; the manifest records the receipt.

Our position is that the two concerns should be mechanically separated:

\begin{itemize}
\item \textbf{Generation is unprivileged.} Anything --- a template engine, a
  human, an LLM --- may produce a \emph{candidate}. Candidates carry
  provenance and a content hash but no standing.
\item \textbf{Standing is admission.} A candidate acquires standing only by
  passing a fail-closed gate: the Lean kernel accepts it, it contains no
  \texttt{sorry}, and its axiom footprint is pinned to an authorized set.
  Every failure is a typed refusal, enumerated in a closed inductive type.
\item \textbf{Claims are scoped.} A release manifest records, per
  declaration, whether it is \emph{proven} (kernel-admitted, axiom-audited)
  or merely \emph{stated} (the formal statement exists; the proof
  obligation is open). The paper's own evaluation is rendered from the
  manifest, not asserted in prose.
\end{itemize}

We call this discipline \emph{mathematical manufacturing}: the pipeline is
the product as much as the library is. The statement--proof split, the typed
refusals, and the manifest are what make an overnight, largely unattended
run reportable at all.

\paragraph{Scope.} This paper does not claim to formalize all of process
intelligence completely, nor does it claim the manufacturing pipeline is
applicable beyond the case study reported here. It introduces a
mathematical manufacturing pipeline and demonstrates it by generating a
first Lean~4 process-intelligence foundation from an existing typed
implementation inventory. Table~\ref{tab:stages} makes the pipeline
concrete: each stage names the artifact it produces, not a metaphor.

\begin{table}[h]
\centering
\small
\begin{tabular}{ll}
\toprule
Manufacturing stage & Concrete artifact \\
\midrule
Source inventory & typed Rust implementation (wasm4pm-compat) \\
Rendering & \textsf{ggen} SPARQL-in-Tera templates (render Lean modules from the TTL catalog) \\
Admission & Lean module creation (\texttt{lake env lean}) \\
Inspection & \texttt{lake build} \\
Refusal & typed \texttt{Refusal}, negative fixtures \\
Evidence & process trace, release manifest \\
Certification & \texttt{CertifiedRelease} object, axiom audit \\
Publication & generated paper evaluation section \\
\bottomrule
\end{tabular}
\caption{Mathematical manufacturing stages and their artifacts.}
\label{tab:stages}
\end{table}

This paper makes four bounded contributions.

First, it defines \emph{mathematical manufacturing} as a typed admission
discipline for formal artifacts. Generated output, human patches,
templates, and agent proposals are treated uniformly as candidates.
Standing is acquired only through admission: kernel acceptance, no-sorry
census, axiom audit, negative fixtures, manifest completeness, and
replayable receipt (Section~\ref{sec:law}).

Second, it introduces \textsf{mfact}, a Lean~4 framework that makes this
discipline explicit through candidate, refusal, admission, manifest,
certified-release, and valid-objection types (Section~\ref{sec:mfact}).
The framework is intentionally small: it does not attempt to replace the
Lean kernel, but to make the pre-kernel and post-kernel standing surface
explicit.

Third, it introduces \textsf{procint}, the first manufactured domain
produced by \textsf{mfact}: a process-intelligence library covering Petri
nets, workflow nets, event logs, conformance, and OCEL-style
object-centric logs (Section~\ref{sec:procint}). The novelty is not only
the domain coverage --- to our knowledge, based on a search of the Lean,
Coq, and Isabelle library indexes at the time of the run, no prior
mechanization of workflow-net soundness or OCEL exists in these systems
--- it is the coupling of domain formalization with an explicit
manufacturing receipt. A prior hand-written mechanization, if found, would
be adjacent but not equivalent unless it provides the same admission,
refusal, manifest, axiom, fixture, and replay surfaces.

Fourth, it reports a manifest-generated evaluation
(Section~\ref{sec:eval}) closed by a kernel-checked Standing Quadrature
witness. The paper's evaluation numbers are not prose claims; they are
rendered from the release manifest, and a generated Lean module
(\texttt{ProcInt.Release.Quadrature}) is refused by the kernel if the
declaration catalog, the audited corpus, and the manifest ever diverge.
This closes the loop from declaration to admission to receipt to replay,
and makes the paper itself an artifact of the manufacturing pipeline it
describes.

\section{The Receipted Manufacturing Law}
\label{sec:law}

The governing law of the pipeline is
\[
R_B \vdash A = \mu(O^*_B)
\]
where \(B\) is the declared boundary (the Lean packages, the declaration
catalog, the manifest schema, the source inventory, the audited theorem
set); \(O\) is raw observation --- candidate output from a human, an LLM,
a template, a repository, or an ontology fragment; \(O^*_B\) is admitted
observation inside \(B\): typed, checked, bounded, replayable; \(\mu\) is
the lawful manufacturing transform (\textsf{ggen} rendering, Lean
admission, Lake build, audit); \(A\) is an artifact with standing (an
admitted theorem, module, manifest entry, or release); and \(R_B\) is the
receipt (build log, manifest hash, axiom audit, negative fixture, replay
evidence). Raw generated output is \(O\), not \(O^*\). \emph{The pipeline
exists to refuse the gap.}

In this architecture, \textsf{ggen} is not a code generator in the usual
sense. It is the projection half of the manufacturing function \(\mu\):
it maps a declared obligation graph into candidate Lean surfaces. Lean
supplies the admission half. Projection without admission is not
standing; admission without projection does not scale.

\paragraph{Rice boundary and manufactured decidability.}
Rice's theorem bars nontrivial semantic decision procedures over
arbitrary programs. The manufacturing response is not to evade this
boundary, but to change the object of judgment. \textsf{mfact} does not
decide arbitrary semantic truth of arbitrary generated code. It decides
whether a bounded candidate has been admitted by a specific kernel, under
a declared axiom policy, with a complete manifest receipt. The
undecidable question ``is this arbitrary generated artifact
mathematically correct?'' is replaced by the decidable question ``did
this candidate pass the declared admission gate inside boundary \(B\)?''
Undecidability is not defeated; it is fenced, and bounded standing is
manufactured inside the fence.

\section{Related Work}

\paragraph{Formal libraries and their pipelines.}
Mathlib~\cite{mathlib2020} is the reference point for large-scale
collaborative formalization in Lean; its quality regime is human review plus
continuous integration. CSLib~\cite{cslib2026} extends the Lean ecosystem
toward computer science, including labeled transition systems, traces, and
bisimulation, which \textsf{procint} builds on for its semantics layer. Lean~4
itself is described by de Moura and Ullrich~\cite{moura2021lean4}. The
closest neighbor to \textsf{mfact} is the agent pipeline of Meek et
al.~\cite{meek2026formalizing}, which formalizes numerical analysis with
coding agents and audits quality beyond kernel acceptance. We differ in
mechanism rather than diagnosis: where their audit is a post-hoc quality
assessment over agent output, \textsf{mfact} is a \emph{typed pipeline} ---
candidates, refusals, and admissions are Lean data types, the statement
corpus lives in an RDF catalog curated by LLM-assisted lanes (each
candidate kernel-verified before being recorded, rather than accepted on
the strength of agent output alone), and the proven/stated split is
enforced by the release gate rather than estimated afterward. The two
approaches are complementary: their semantic audits address exactly the
residual risk our pipeline does not eliminate, namely that a curated
statement, though kernel-admitted, is itself a faithless or vacuous
rendering of the intended theorem.

\paragraph{Process intelligence.}
The mathematical canon \textsf{procint} formalizes is due to a small number
of primary sources. Petri-net structure theory --- state equation,
invariants, boundedness, liveness --- follows Murata's
survey~\cite{murata1989}. Workflow nets and their soundness are van der
Aalst's~\cite{vanderaalst1997}; the decidability and classification
landscape is surveyed in~\cite{vanderaalst2011soundness}, and the process
mining setting in~\cite{vanderaalst2011pm}. Alignment-based conformance
checking follows Carmona et al.~\cite{carmona2018}. Object-centric event
data follows the OCEL~2.0 specification~\cite{ocel2} and object-centric
Petri nets~\cite{ocpn2020}; partially ordered workflow models follow
POWL~\cite{powl2023}. None of these, to the best of our search, had a prior
mechanization in a proof assistant.

\section{The mfact Framework}
\label{sec:mfact}

\textsf{mfact} is deliberately small: seven modules and a CLI. Its job is to
make the pipeline's epistemics first-class Lean objects.

\paragraph{Candidates and refusals.}
A \texttt{Candidate} is content plus provenance (producer identity, content
hash); producers are untrusted by construction. \texttt{Refusal} is a closed
inductive type --- \texttt{kernelRejected}, \texttt{sorryPresent},
\texttt{unauthorizedAxiom}, \texttt{missingEvidence},
\texttt{malformedModel}, \texttt{invalidFixture} --- so a failed admission is
always a value that can be logged, hashed, and re-queued, never an untyped
error string. There is no default case: an unrecognized failure mode is
itself a refusal to classify, not a silent pass.

\paragraph{Admission with a covers proof.}
An \texttt{Admission} bundles the evidence for one candidate: the kernel
acceptance record, the sorry census, the axiom audit output, and a
\texttt{covers} field --- a proof term witnessing that the evidence entries
cover the obligations the gate imposes, so an admission cannot be assembled
with a missing check. The gate function \texttt{admit} is total: it returns
either an \texttt{Admission} or a \texttt{Refusal}, with no third outcome.

\paragraph{Manifest, release, objections.}
A \texttt{Manifest} (JSON-serializable via \texttt{ToJson}) lists artifacts
with hashes, axiom sets, fixture results, and the proven/stated status of
each declaration. A \texttt{CertifiedRelease} pairs a manifest with its
audit obligations. \texttt{ValidObjection R} is a closed inductive family
whose constructors are the only admissible objection forms against a
release \texttt{R} (an unaudited theorem, a missing hash, an unauthorized
axiom, a failing fixture); each constructor demands evidence. The
per-release theorem
\begin{center}
\texttt{theorem no\_valid\_objection : IsEmpty (ValidObjection R)}
\end{center}
is discharged by refuting each constructor against the corresponding checked
field of \texttt{R}. For the release described in this paper the theorem is
proven and, per \texttt{\#print axioms}, depends on no axioms at all --- not
even the three standard Lean/Mathlib axioms. The theorem's scope is exactly
its constructors: it says the \emph{enumerated} objection forms are refuted,
not that the release is beyond all criticism.

\section{procint: the Mathematical Canon}
\label{sec:procint}

In this paper, process intelligence does not mean dashboarding over event
logs. It means a mathematical substrate for deciding, replaying, and
explaining process standing: what happened, what was admitted, what was
refused, which object relations were involved, which conformance
obligations passed, and which receipts authorize the resulting claim.
\textsf{procint} depends on Mathlib and CSLib and is organized as follows
(module map abridged):

\begin{center}
\small
\begin{tabular}{ll}
\toprule
Namespace & Content \\
\midrule
\texttt{ProcInt.Foundations} & metrics as \(\{q : \mathbb{Q} \mathrel{/\!/} 0 \le q \land q \le 1\}\), identifiers, LTS \\
\texttt{ProcInt.Petri} & nets, firing, reachability, state equation, invariants, \\
 & boundedness, liveness, OCPN, stochastic nets, unfoldings \\
\texttt{ProcInt.Workflow} & WF-nets, short-circuit construction, soundness \\
\texttt{ProcInt.Logs} & events, traces (monotone timestamps), event logs, XES \\
\texttt{ProcInt.Models} & process trees, POWL, DFGs, causal nets, BPMN, Declare \\
\texttt{ProcInt.Conformance} & moves, alignments, cost, fitness, token replay \\
\texttt{ProcInt.Ocel} & OCEL~2.0, E2O/O2O relations, flattening, lifecycles \\
\texttt{ProcInt.Ocpq} & object-centric process queries \\
\texttt{ProcInt.Analytics} & temporal, causality, correlation, cubes, streaming \\
\texttt{ProcInt.Registry} & algorithm and breed registries (cross-product tables) \\
\bottomrule
\end{tabular}
\end{center}

\paragraph{Petri nets.}
A net is places, transitions, and weighted flow over finite types; a marking
is a finitely supported function \(M : P \to_{f} \mathbb{N}\), which unifies
the weighted and (via richer codomains) colored cases and gives the marking
algebra of Murata~\cite{murata1989} --- firing, the state equation
\(M' = M + A^{\mathsf T} u\), and place/transition invariants --- as
\(\mathsf{Finsupp}\) lemmas with truncated subtraction.

\paragraph{Workflow nets and soundness.}
A \texttt{WfNet} carries a Petri net as a field (composition, not
extension) together with the source-place, sink-place, and connectivity
conditions of van der Aalst's Definition~6--7~\cite{vanderaalst1997}.
Following the observation that the classical soundness clauses are logically
independent, soundness is \emph{not} a single conjunction-shaped definition
but three separately named predicates --- option to complete, proper
completion, and no dead transitions --- with \texttt{Sound} as their
explicit conjunction. This keeps partial results honest: a lemma about one
clause is stated about that clause. The crown-jewel target of the run is the
short-circuit characterization (soundness of \(N\) iff liveness and
boundedness of the short-circuited net, \cite{vanderaalst1997}, Lemma~8);
its per-direction status at release time is recorded in the manifest and
reported in Section~\ref{sec:eval}, not asserted here.

\paragraph{Conformance.}
Alignments follow Carmona et al.~\cite{carmona2018}: moves are synchronous,
log-only, model-only, or silent; costs induce optimal alignments; fitness is
a \(\mathsf{Between01}\) value by construction, so boundedness of the metric
is a type-level fact rather than a theorem obligation.

\paragraph{OCEL 2.0.}
Object-centric event logs follow the OCEL~2.0
specification~\cite{ocel2} (Definitions~1--2): typed events and objects,
event-to-object and object-to-object relations with qualifiers, and
attribute maps; flattening onto a classical log and object lifecycles are
derived operations. Object-centric Petri nets follow~\cite{ocpn2020}.

\paragraph{Canon status.} Table~\ref{tab:canon} states, per area, whether
\textsf{procint} v1 formalizes it, formalizes a seed instance of it, or
leaves it to future work. The claim of this paper is a first combinatorial
process-intelligence foundation, not a complete formalization of process
intelligence; per-declaration status (proven/stated) within each formalized
area is in the generated evaluation (Section~\ref{sec:eval}).

\begin{table}[h]
\centering
\small
\begin{tabular}{ll}
\toprule
Area & v1 status \\
\midrule
Transition systems, traces, reachability & formalized (Foundations, Petri) \\
Petri nets (marking, firing, invariants) & formalized \\
Workflow nets and soundness & formalized; short-circuit theorem partial (Sec.~\ref{sec:eval}) \\
Alignment-based conformance & formalized (moves, cost, fitness) \\
OCEL~2.0 event/object relations & formalized \\
Object-centric Petri nets & seed instance \\
Process trees, POWL, DFGs, BPMN, Declare & seed instances \\
Object-centric process queries (OCPQ) & seed instance \\
Temporal/causal/streaming analytics & seed instances \\
\bottomrule
\end{tabular}
\caption{Process-intelligence canon coverage in \textsf{procint} v1.}
\label{tab:canon}
\end{table}

\section{The Manufacturing Run}
\label{sec:run}

The run that produced \textsf{procint} has five moving parts.

\paragraph{The TTL catalog as canonical source.}
The canonical source of \textsf{procint} is not a tree of \texttt{.lean}
files but an RDF/Turtle declaration catalog (\texttt{procint:} vocabulary):
one record per declaration, carrying its Lean body (\texttt{leanCode}),
its module, its proof status, its axiom-audit expectation, and citations
back to the primary literature. Declaration bodies enter the catalog by
curation --- LLM-assisted lanes draft them and each candidate is
kernel-tested before being recorded --- so the catalog holds verified
candidates, not free-form generation. The cross-product surfaces
(registries of algorithms and model breeds) are fully data-driven rows of
the same graph.

\paragraph{SPARQL-in-Tera rendering.}
The \textsf{ggen} engine renders Lean modules from templates whose
frontmatter carries SPARQL \texttt{SELECT} queries against the catalog and
whose bodies are Tera templates over the result bindings. One template
yields one \texttt{.lean} file per \texttt{procint:Module}; separate
templates render the root import index, the registry tables, and the
axiom-audit file --- for every audited theorem, a \texttt{\#guard\_msgs in
\#print axioms} block pinned to \texttt{[propext, Classical.choice,
Quot.sound]}, so any transitive \texttt{sorryAx} or unauthorized axiom
breaks the build. The rendered corpus is a \emph{candidate artifact}: it is
not trusted because it was rendered, and \textsf{ggen} certifies nothing.
It acquires standing only through kernel admission. In one line:
\textsf{ggen} renders, Lean admits, \textsf{mfact} certifies.

\paragraph{LLM lanes as untrusted producers.}
Module families (Petri core, workflow, logs and conformance, OCEL,
models, analytics, fixtures) run as independent lanes with no barriers
between them. Each lane's agent hand-drafts the actual structures,
definitions, and proofs for its module family, but writes nothing directly
into the library: every piece of Lean code it proposes is first
kernel-tested via \texttt{lake env lean} on a scratch file, and only a
declaration that passes this check is recorded --- as a \texttt{leanCode}
string literal, together with its proof status --- in the lane's catalog
fragment. The catalog records what the lane authored and what the kernel
confirmed; it authors nothing itself. Statements the lane cannot prove ship
as \texttt{sorry}-bearing stubs marked \emph{stated} in the fragment.

\paragraph{Admission and the repair loop.}
An integrator loop repeatedly merges fragments, re-renders, and runs
\texttt{lake build}. Build errors are parsed and returned to the owning
lane as repair context, with bounded retries per stub and per lane;
exhausted stubs are downgraded --- the statement stays in the library, the
claim drops to \emph{stated}, and the declaration leaves the audited set.
Nothing is deleted and nothing passes silently: the terminal state of every
stub is either kernel-admitted or an explicit downgrade receipt.

\paragraph{Certification.}
The release gate rebuilds both packages, builds the axiom-audit files,
runs the negative fixtures (\texttt{\#guard\_msgs (error)} blocks, the Lean
analog of compile-fail tests), and emits the release manifest with content
hashes. A negative control --- injecting a \texttt{sorry} into a scratch
copy of an audited theorem --- must cause the gate to refuse; a gate that
cannot fail certifies nothing.

\section{Use of Generative AI Tools}
\label{sec:ai-disclosure}

Claude Code and related language-model tooling were used to generate
candidate patches, proof sketches, ontology fragments, documentation
drafts, and release metadata during the run described in
Section~\ref{sec:run}. These outputs were treated throughout as untrusted
candidate artifacts, per the separation in Section~\ref{sec:mfact}: they
were admitted into \textsf{procint} only when accepted by the Lean kernel,
the sorry census, the axiom audit, schema validation against the
\texttt{procint:} ontology, and the negative-fixture and release-gate
checks of Section~\ref{sec:run}. No language-model output is part of the
trusted base (Table~\ref{tab:stages}); the trusted base is the Lean kernel,
Mathlib and CSLib at their pinned revisions, and the \texttt{lake}/\texttt{ggen}
toolchain. The author reviewed and accepts responsibility for all submitted
content, code, citations, and claims in this paper and in the accompanying
repository.

\section{Evaluation}
\label{sec:eval}

Every number in this section is rendered from the release manifest
(\texttt{release-manifest.json}) of the run; none is written by hand. The
generation pipeline fails closed: if the manifest is absent or stale, this
section does not render and the paper does not build for release.

\ifreleasebuild
  \input{generated/evaluation}
\else
  \InputIfFileExists{generated/evaluation}{}{%
  \begin{center}\fbox{\itshape Draft build: generated/evaluation.tex missing.}\end{center}}
\fi

\paragraph{Final status.} Table below is generated from the release
manifest and standing records; no cell is hand-authored.
\ifreleasebuild
  \input{generated/final_status}
\else
  \InputIfFileExists{generated/final_status}{}{%
  \begin{center}\fbox{\itshape Draft build: generated/final\_status.tex missing.}\end{center}}
\fi

\subsection{Standing Quadrature}
\label{sec:quadrature}

The release is not evaluated only by whether Lean builds. It is evaluated
by whether the declaration catalog, admitted Lean corpus, process
evidence, release manifest, and paper claims form a \emph{closed standing
graph}. The quadrature check rejects orphan declarations, unsupported
claims, untraced artifacts, and manifest--paper divergence. Closure is
itself kernel-checked: a generated witness module
(\texttt{ProcInt.Release.Quadrature}) carries the name-sorted proven
surfaces of the catalog, the manifest, and the axiom audit as literal
lists, and proves their pairwise equality by \texttt{rfl} --- any orphan
on any side makes the lists differ and the kernel refuses the witness.
The same module encodes the manufacturing run as a \textsf{procint}
process trace and proves, with the corpus's own \texttt{DeclareConstraint}
semantics, that the run conforms to its declared process model:
\textsf{procint} models the process by which \textsf{mfact} manufactured
\textsf{procint}. Table~\ref{tab:quadrature} is generated from
\texttt{quadrature.json}; three negative controls (a corrupted evaluation
number, a claim with no evidence, a catalog declaration with no rendered
symbol) each cause a typed refusal.

\ifreleasebuild
  % RENDERED by ggen sync from the quadrature graph (quadrature-pack).
% Do not edit: regenerate via `just standing-quadrature`.
\begin{table}[h]
\centering
\small
\begin{tabular}{lrl}
\toprule
Surface pair & Count & Result \\
\midrule
TTL to Lean & 145 & PASS \\
Lean to Audit & 145 & PASS \\
Audit to Manifest & 145 & PASS \\
Manifest to Paper & 6 & PASS \\
Artifact to Process Event & 193 & PASS \\
Claim to Evidence & 5 & PASS \\
\bottomrule
\end{tabular}
\caption{Standing Quadrature: PASS over
145 proven declarations, 5 traced paper
claims, and 6 evaluation numbers. The closure is
kernel-checked by the rendered witness
\texttt{ProcInt.Release.Quadrature}.}
\label{tab:quadrature}
\end{table}

\else
  \InputIfFileExists{generated/quadrature}{}{%
  \begin{center}\itshape Draft build: generated/quadrature.tex missing.\end{center}}
\fi

\section{Falsifier and Valid Objection Surface}
\label{sec:falsifier}

A valid objection to this release must show at least one of:
(1) a declared admitted theorem contains \texttt{sorry};
(2) a theorem depends on unauthorized axioms;
(3) a manifest hash does not match the artifact;
(4) a negative fixture passes when it should fail;
(5) a generated evaluation number does not match the manifest;
(6) a stated declaration is described as proven;
(7) the catalog-to-Lean rendering changes the intended theorem statement;
(8) the release cannot be replayed under pinned dependencies.
A local objection to one proof, citation, or template is not
automatically a refutation of mathematical manufacturing. It becomes a
refutation only if it crosses the admission boundary and invalidates a
claimed receipt.

\section{Limitations and Standing}
\label{sec:limitations}

\paragraph{Proven versus stated.}
The library's standing is exactly the manifest's proven/stated split; prose
in this paper adds nothing to it. \emph{Stated} declarations are formal
statements whose proof obligations are open --- they compile, they are
citable as definitions of the intended property, and they carry no more
authority than that.

\paragraph{Crown-jewel scope.}
The short-circuit soundness characterization is reported per direction. If
only one direction is admitted at release time, the claim is that direction
and nothing more; the converse remains a stated obligation in the manifest.
\ifreleasebuild
  % GENERATED FILE — DO NOT EDIT BY HAND (source: quadrature graph via ProcInt.crownJewel_status)
The release-time classification of the short-circuit soundness
characterization is \textbf{\CrownJewelStatus} (ground truth:
\texttt{ProcInt.crownJewel\_status}); this fragment is generated from the
catalog and cannot report it as proven while the obligation is open.

\else
  \InputIfFileExists{generated/crown_jewel_status}{}{}
\fi

\paragraph{What admission does not check.}
Kernel admission plus the axiom audit establishes that admitted terms prove
the stated theorems under the pinned axioms. It does not establish that
the catalog-curated statements faithfully render the informal canon; that
correspondence rests on citation-annotated statement curation against the
primary sources and is exactly the kind of residual risk the semantic
audits of~\cite{meek2026formalizing} target. Fidelity of statements is a
review obligation, not a theorem.

\paragraph{Scope of the novelty claim.}
``First mechanization'' claims are scoped to a search of the Lean
(Mathlib/CSLib), Coq, and Isabelle (AFP) ecosystems at run time; we claim a
negative search result, not a proof of absence.

\paragraph{Exclusions.}
This paper does \emph{not} claim any of the following.

\begin{center}
\small
\begin{tabular}{ll}
\toprule
Exclusion & Reason \\
\midrule
LLMs prove mathematics autonomously & LLMs are untrusted producers \\
\textsf{procint} formalizes all of process intelligence & v1 boundary is manifest-scoped \\
Kernel admission proves informal intent & statement fidelity is a review obligation \\
Rice's theorem is violated & arbitrary semantic correctness is not decided \\
All objections are impossible & only enumerated release objections are refuted \\
\textsf{ggen} output has standing by generation & only admitted output has standing \\
A stated theorem is proven & stated means the obligation remains open \\
\bottomrule
\end{tabular}
\end{center}

\section{Availability}
\label{sec:availability}

The \textsf{mfact} and \textsf{procint} packages, the ontology, the
templates, the release manifest, and the source of this paper (including
the generated evaluation) are available in the \texttt{mfact} repository,
pinned to Lean~4 v4.31.0 with Mathlib and CSLib at fixed revisions. The
negative-control script that demonstrates the gate refusing a
\texttt{sorry}-injected artifact ships with the release.

\paragraph{Reproducibility appendix.}
\begin{itemize}
\item \textbf{Lean:} \texttt{leanprover/lean4:v4.31.0}.
\item \textbf{Mathlib:} \texttt{leanprover-community/mathlib4} at the commit
  tagged \texttt{v4.31.0}.
\item \textbf{CSLib:} \texttt{leanprover/cslib}, pinned to its last commit
  on the \texttt{v4.31.0} toolchain.
\item \textbf{Build:} \texttt{lake build} in \texttt{mfact/} and
  \texttt{procint/}.
\item \textbf{Rendering:} \texttt{ggen sync run} against
  \texttt{packs/lean-math-pack} and the merged \texttt{procint:} catalog.
  This renders the candidate Lean corpus (modules, index, axiom audit,
  registry tables) from the TTL catalog; the rendered corpus is admitted
  only by the subsequent \texttt{lake build}.
\item \textbf{Certification:} \texttt{lake exe mfact certify
  release-manifest.json gates.json}.
\item \textbf{Artifact and manifest hashes:} recorded per-declaration
  (BLAKE3) and as a genesis-folded release hash in
  \texttt{release-manifest.json}; reproduced in the generated evaluation
  (Section~\ref{sec:eval}).
\item \textbf{Expected output:} \texttt{CERTIFIED\_RELEASE=PASS} with
  \texttt{SORRY\_COUNT=0} and \texttt{AXIOM\_AUDIT=PASS}; a negative control
  with an injected \texttt{sorry} must instead exit refused.
\end{itemize}

\ifreleasebuild
  \input{generated/availability}
\else
  \InputIfFileExists{generated/availability}{}{}
\fi

\section{Conclusion}

\ifreleasebuild
  \input{generated/conclusion}
\else
  \InputIfFileExists{generated/conclusion}{}{%
  \begin{center}\fbox{\itshape Draft build: generated/conclusion.tex missing.}\end{center}}
\fi

\bibliographystyle{plain}
\bibliography{refs}

\end{document}
